\label{cha:DistributedComparisonFunctions}
\subsection{Comparison Functions}
\label{sub:ComparisonFunctions}
\textit{Comparison Functions} act like \textit{Point Functions} (\ref{sub:PointFunctions}) just for intervals.
This means:
\[
f^{<}_{\alpha}(i) = \begin{cases} 
      1 & \text{if } i < \alpha, \\
      0 & \text{else}.
   \end{cases}
\]

Recall the binary tree implementation of \textit{Point Functions}.
Please note, the following implementation only works if the leaves are placed in ascending order.
This means, the smallest leaf is on the left while the biggest leaf is on the right.
In our implementation, the leafs represent indices with the leaf on the left representing index $0$ and the rightmost leaf representing index $2^n$ (assuming a tree of depth $n$).

\subsection{Distributed Comparison Functions (DCF)}
\label{sub:DistributedComparisonFunctions}
The \textit{FSS} (\ref{ch_FunctionSecretSharing}) for secretly sharing \textit{Comparison Functions} is called \textit{Distributed Comparison Function (DCF)}.
On a high level, both shares are represented by binary trees.
If one evaluates the function, she traverses the tree as described in \ref{sub:TraversingTheDPFTree}.
Using \textit{Point Functions} (\ref{sub:PointFunctions}), one is interested in the flag bits, correction words (including the final correction word), and the label of the leaf.
Now, one is interested in whether she stays on the path or not while traversing the tree.
Since the leafs are placed in ascending order, leaving the path to leaf $\alpha$ to the left implicates that the given value is smaller than threshold $\alpha$.
Thus, the function evaluates to $1$.
Leaving the path to $\alpha$ to the right or staying on the path to alpha with a given input $i$ implicates $i\geq\alpha$.
The function evaluates to $0$.

\subsection{DCF Construction}
\label{sub:CDPFConstruction}
Multiple \textit{DCF} constructions coexist.
This research focuses on the construction by Kyle et al. \cite{kyle2023grotto}.
Goal of this construction is to find out whether a shared secret $x$ is greater or less than zero.
The benevolent reader might also be interested in the more prominent \textit{DCF} construction of Boyle et al. \cite{boyle2021function}.

Kyle et al. \cite{kyle2023grotto} are more favorable in their communication cost and ease of implementation.
Therefore, this construction is used in the implementation.
For this purpose one uses a pair of the already mentioned \textit{DPFs} (\ref{sub:DistributedPointFunctions}), and expands them to meet the functional requirements of \textit{DCFs}.
Also, calculating the result only requires one single word to be exchanged.

First, both \textit{DPFs} are generated on the same random target value.
For ease of explanation, we will take a look at a construction using a single \textit{DPF} first and evaluate on 64-bit signed integers only.
Both parties $P_0 \text{ and }P_1$ hold a share of the random target value $target_0\text{ and } target_1$ with $target = target_0+target_1$.
Given a secret-shared input $x$ with $x=x_0+x_1$ each party calculates $target_i-x_i$.
Both parties exchange this value.
This allows them to reconstruct $S=target-x=(target_0-x_0)+(target_1-x_1)$.

To simplify the explanation consider a bit-vector $V$ with length $2^{64}$.
It is populated with 1s in the range of $[S,S+2^{63}-1]$ and 0s everywhere else.
This corresponds to the partitioning of the two's complement for a 64-bit signed integer.
Let us assume the result of the \textit{DPF} to be the dot-product of this vector $V$ and the \textit{DPF}.
This means, 0 everywhere but the target-index of the \textit{DPF}.

If one now evaluates the \textit{DPF} on $target$, which is equal of calculating $V[target]$, one receives a value of either $0$ or $1$.
Assume $V[target]=1$.
This means $target = S+k$ for some $k\in [1,2^{63}-1]$.
We conclude that $x>0$ since $S$ is shifted to the left and therefore $V$ evaluates to $0$.
If $V[target]=0$ we conclude $x<=0$ since $S=target-(-x) = target + x$ which shifts the target into the $0$ area of $V$.